% Sección 4: Metodología analítica y diseño de validación temporal
% evidencia:E08,E09,E11,E12,E13,E14

\section{Metodología analítica y diseño de validación temporal}

El diseño predictivo de MediFlow separa estrictamente tres familias de modelos que cumplen propósitos distintos dentro del sistema: una escalera experimental puntual basada en modelos de conteo (Poisson) para contrastar hipótesis y seleccionar variables externas; un conjunto de modelos cuantílicos servido en la aplicación web para proveer intervalos y descomposiciones explicativas; y un modelo independiente sin rezagos cortos para proyectar horizontes de hasta cuatro meses. Todas las familias comparten la prohibición de particiones aleatorias y se evalúan mediante validación temporal de origen móvil (\textit{walk-forward}) \parencite{hyndman_fpp3}.

\subsection{Validación temporal de origen móvil y política anti-fuga}

Sea $\{y_{s,t}\}$ la serie regular diaria de llegadas únicas del servicio $s$ en la fecha $t$. Para un pliegue $k$ definido por un origen temporal $o_k$ y un horizonte $h$, la partición temporal en \texttt{src/hun/models/validacion.py} impone:
\begin{equation}
\label{eq:walk_forward}
\mathcal{D}_{\text{train}}^{(k)} = \{(s, t) : t < o_k\}, \qquad
\mathcal{D}_{\text{test}}^{(k)} = \{(s, t) : o_k \le t < o_k + h\},
\end{equation}
de modo que ningún registro de entrenamiento incluye fechas iguales o posteriores al origen del pliegue. En la evaluación operativa a corto plazo se utilizan \textbf{26 pliegues} con orígenes comprendidos entre \texttt{2025-07-25} y \texttt{2026-07-10}, paso fijo de 14 días y horizonte $h = 7$ días \parencite{mediflow2026demanda}. Debido a que el paso de 14 días es múltiplo de siete, todos los orígenes caen en viernes; en consecuencia, el día 1 del horizonte corresponde siempre a viernes, el día 2 a sábado, el día 3 a domingo y así sucesivamente hasta el jueves. El día del horizonte y el día de la semana quedan estructuralmente confundidos, por lo que las variaciones de error entre los días 1 y 7 no deben atribuirse únicamente al alejamiento temporal del pronóstico. En el horizonte extendido de largo plazo se emplean \textbf{13 pliegues} entre \texttt{2025-03-24} y \texttt{2026-03-19} con paso de 30 días y horizonte $h = 120$ días; como el horizonte cuadruplica el paso, las ventanas de prueba consecutivas se solapan extensamente y no constituyen muestras independientes \parencite{mediflow2026largoplazo}.

La política anti-fuga implementada en \texttt{construir\_features\_demanda} de \texttt{src/hun/models/demanda.py} enmascara explícitamente el vector objetivo $y_{s,t}$ como faltante para todo $t \ge o_k$ antes de calcular desplazamientos semanales (\texttt{llegadas\_lag\_7} hasta \texttt{llegadas\_lag\_56}, además de \texttt{rezago\_prom\_4s} y \texttt{rezago\_mediana\_8s}). Así, para predecir cualquier día dentro de $[o_k, o_k + 7\,\text{d})$, los rezagos mínimos de siete días consultan exclusivamente observaciones estrictamente anteriores al origen. Las series epidemiológicas de SIVIGILA incorporan rezagos de dos y tres semanas y el modelo de largo plazo sustituye cualquier observación meteorológica posterior por medianas climáticas por día del año calculadas con fechas anteriores a cada corte. Como límite explícito, este mecanismo protege la evaluación \textit{walk-forward} de 1--7 días, pero no hereda automáticamente esa garantía cuando la aplicación web realiza inferencia puntual en fechas posteriores al corte usando medianas históricas de respaldo (\texttt{medianas\_dw}) ni cuando el modelo final entrenado con toda la historia se consulta sobre fechas ya vistas \parencite{mediflow2026bundle,mediflow2026informe}.

\subsection{Métricas de evaluación y techo Poisson idealizado}

Dado que las series hospitalarias exhiben días con volumen nulo o muy bajo (por ejemplo, Rehabilitación en domingos), las métricas basadas en porcentajes individuales como MAPE divergen. Siguiendo la decisión técnica \texttt{DEC-003}, la evaluación adopta el Error Porcentual Absoluto Ponderado (WAPE), acompañado de MAE, RMSE, sesgo y cobertura empírica:
{\singlespacing
\begin{equation}
\label{eq:metricas}
\begin{aligned}
\text{WAPE} &= \frac{\sum_{i=1}^{n} |y_i - \hat{y}_i|}{\sum_{i=1}^{n} |y_i|}, &
\text{Exactitud} &= 1 - \text{WAPE}, \\
\text{MAE} &= \frac{1}{n}\sum_{i=1}^{n} |y_i - \hat{y}_i|, &
\text{RMSE} &= \sqrt{\frac{1}{n}\sum_{i=1}^{n} (y_i - \hat{y}_i)^2}, \\
\text{Sesgo} &= \frac{1}{n}\sum_{i=1}^{n} (\hat{y}_i - y_i), &
\text{Cobertura} &= \frac{1}{n}\sum_{i=1}^{n} \mathbf{1}\!\left\{\hat{q}_{0{,}10,i} \le y_i \le \hat{q}_{0{,}90,i}\right\}.
\end{aligned}
\end{equation}
}
En este proyecto, «exactitud» significa estrictamente $1 - \text{WAPE}$ y no porcentaje de días acertados individualmente; un sesgo positivo indica sobreestimación promedio y la meta interna de 80\,\% es un criterio interno de trabajo del equipo, no un umbral de la rúbrica oficial \parencite{mediflow2026demanda,hun2026terminos}. La serie residual \texttt{Otros} (menos de una llegada diaria) se excluye de las métricas globales para no distorsionar los promedios agregados.

Para contextualizar servicios de bajo volumen como Tratamiento Urgencias ($\approx 15{,}2$ llegadas/día), \texttt{techo\_exactitud\_poisson} calcula un techo idealizado suponiendo conteos independientes $X \sim \text{Poisson}(\mu)$ con tasa verdadera conocida $\mu = \bar{y}\cdot d$ ($d=1$ diario, $d=7$ semanal):
\begin{equation}
\label{eq:techo_poisson}
\text{Techo}(\mu) = 1 - \frac{\mathbb{E}\bigl[|X - \mu|\bigr]}{\mu}.
\end{equation}
Puesto que el cálculo presupone conocer la tasa verdadera exacta y que la única fuente de variabilidad es el ruido puro de Poisson sin cambios semanales de nivel, este techo constituye una referencia teórica optimista y no una cota estadística universal \parencite{mediflow2026demanda}.

\subsection{Familia experimental puntual (escalera Poisson y LightGBM)}

La escalera experimental implementada en \texttt{src/hun/models/demanda.py} compara cuatro escalones sobre los mismos 26 pliegues: (1) \texttt{linea\_base}, definida como la mediana del mismo día de la semana durante las últimas ocho semanas; (2) \texttt{poisson\_calendario}, regresión lineal generalizada con enlace logarítmico (\texttt{PoissonRegressor}, $\alpha = 10^{-4}$) ajustada \textbf{de forma estratificada por servicio} (evitando mezclar los cierres festivos de Rehabilitación con la operación 24/7 de Urgencias); (3) \texttt{lgbm\_todos\_externos}, basada en LightGBM \parencite{ke2017lightgbm} con pérdida Poisson (\texttt{n\_estimators=200}, \texttt{learning\_rate=0.05}, \texttt{num\_leaves=15}, \texttt{min\_child\_samples=15}); y (4) \texttt{lgbm\_final}, que restringe las variables externas a los grupos aprobados por el protocolo.

Dos componentes algorítmicos distinguen a los modelos LightGBM de la escalera puntual frente a una configuración genérica. Primero, la ponderación exponencial por recencia asigna a cada fila de entrenamiento un peso $w_i = 0{,}5^{\text{edad}_i / 180}$ (vida media de 180 días). Segundo, la función \texttt{ajustar\_predecir\_lgbm\_nivel} introduce una corrección de nivel local mediante \texttt{init\_score} $= \log(\tilde{y}_{\text{local}})$, donde $\tilde{y}_{\text{local}} = \max(\text{rezago\_mediana\_8s},\, 0{,}5)$, de manera que los árboles predicen en escala logarítmica el factor multiplicativo sobre el nivel reciente de cada servicio y escalan proporcionalmente ante cambios de régimen.

\subsection{Familia cuantílica servida en la aplicación web}

El paquete desplegado para la API web (\texttt{models/modelos\_web\_bundle.joblib}, versión \texttt{2.0}) no sirve el modelo \texttt{lgbm\_final} con pérdida Poisson ni comparte idéntica arquitectura de entrenamiento. En su lugar, \texttt{src/hun/models/explicacion.py} ajusta tres regresores LightGBM independientes mediante pérdida \textit{pinball} o de regresión cuantílica \parencite{koenker1978quantiles,ke2017lightgbm} para los niveles $\tau \in \{0{,}10,\, 0{,}50,\, 0{,}90\}$:
\begin{equation}
\label{eq:pinball}
\rho_\tau(u) = u\bigl(\tau - \mathbf{1}\{u < 0\}\bigr), \qquad u = y - \hat{q}_\tau(x).
\end{equation}
Para evitar cruces entre cuantiles estimados por árboles separados y valores negativos de conteo, se aplica un ordenamiento monotónico y truncamiento posterior:
\begin{equation}
\label{eq:orden_cuantiles}
\bigl(\tilde{q}_{0{,}10},\, \tilde{q}_{0{,}50},\, \tilde{q}_{0{,}90}\bigr) = \max\!\Bigl(0,\; \operatorname{sort}\bigl(\hat{q}_{0{,}10},\, \hat{q}_{0{,}50},\, \hat{q}_{0{,}90}\bigr)\Bigr).
\end{equation}
Este ajuste garantiza coherencia ordinal en las respuestas de la interfaz, pero no certifica ni corrige la calibración probabilística de la cobertura empírica. Para interpretar las salidas, el modelo mediano $\hat{q}_{0{,}50}$ se descompone mediante valores SHAP exactos de árboles (\texttt{TreeExplainer}) \parencite{lundberg2017shap}, expresados como suma aditiva respecto de un valor base $\phi_0$:
\begin{equation}
\label{eq:shap}
f(x) = \phi_0 + \sum_{j=1}^{M} \phi_j(x),
\end{equation}
agrupando las contribuciones en factores legibles para la interfaz. Las atribuciones SHAP describen contribuciones marginales al pronóstico del modelo y no miden efectos causales de intervenciones operativas.

La inspección directa del código y de la lista de columnas \texttt{cols} almacenada en el artefacto \texttt{v2.0} confirma además dos diferencias estructurales frente a \texttt{lgbm\_final}: la función \texttt{entrenar\_modelos\_cuantiles} no emplea desplazamiento logarítmico de nivel local (\texttt{init\_score}) ni pesos exponenciales por recencia, y las 21 columnas efectivamente empaquetadas en el bundle vigente incluyen 5 variables de calendario, 10 variables de rezago y únicamente las 6 variables de \texttt{festivos\_y\_puentes}, omitiendo las 5 variables de \texttt{carnaval\_y\_semana\_santa} \parencite{mediflow2026bundle}. Por ello, sus métricas de desempeño deben reportarse de forma completamente separada del experimento Poisson.

\subsection{Modelo independiente de horizonte extendido}

Para proyectar de agosto a noviembre de 2026 más allá de los últimos rezagos semanales disponibles en la base (\texttt{2026-07-16}), el módulo \texttt{src/hun/models/largo\_plazo.py} entrena un modelo independiente basado en LightGBM Poisson con decaimiento exponencial por recencia (vida media de 180 días, \texttt{min\_child\_samples=20}) sin ningún rezago corto de llegadas. Sus predictores combinan 7 variables de calendario, 6 de \texttt{festivos\_y\_puentes}, la tendencia secular \texttt{tendencia\_anios}, el identificador de servicio y las medianas climáticas históricas por día del año calculadas estrictamente con fechas anteriores a cada corte. Para contener el deterioro de incertidumbre entre uno y cuatro meses, los intervalos parten de regresores cuantílicos base en $p_{05}$ y $p_{95}$ (\texttt{CUANTIL\_BAJO = 0,05} y \texttt{CUANTIL\_ALTO = 0,95}) expandidos alrededor de la estimación puntual $\hat{y}$ según el mes del horizonte $m \in \{1,2,3,4\}$ mediante el factor $1 + 0{,}10\,(m - 1)$, aunque el esquema JSON expuesto a la interfaz rotule los extremos resultantes como \texttt{min\_diario}/\texttt{max\_diario} o intervalo $[p_{10}, p_{90}]$ \parencite{mediflow2026largoplazo}.
